%--------------------------------------------------------------------------------
%
%
%--------------------------------------------------------------------------------
\pagebreak
\unnumberedSection{LDO - Load Offset}

%--------------------------------------------------------------------------------
\begin{iPageDescEntry}{Syntax}
\texttt{LDO RegR, ofs(RegB)} \\
\texttt{LDO RegR, RegX(RegB)}
\end{iPageDescEntry}

%--------------------------------------------------------------------------------
\begin{iPageDescEntry}{Format}
The \texttt{LDO} instruction uses the following instruction formats: \\

\begin{flushleft}
\drawInstrFormat
    {0,1,3,5,6.5,8.5, 16}
    {\OpCodeGrpALU,\OpCodeLDO,RegR,0,RegB, ofs}
\end{flushleft}

\begin{flushleft}
    \drawInstrFormat
        {0,1,3,5,6.5,8.5,9.5,11.5,16}
        {\OpCodeGrpALU,\OpCodeLDO, RegR, 1, RegB, 0, RegX, 0}
    \end{flushleft}
    
\end{iPageDescEntry}

%--------------------------------------------------------------------------------
\begin{iPageDescEntry}{Description}
The \texttt{LDO} instruction loads an offset into a general register. Two formats are available: the immediate-offset format and the register-indexed format. In the immediate-offset format, the signed immediate offset is added to \texttt{RegB} to form the effective address. In the register-indexed format, \texttt{RegX} is added to \texttt{RegB} to form the effective address. The computed value is then stored in \texttt{RegR}.

The addition of base register and offset will wrap around and leave the upper bits of the base register unaffected.

\end{iPageDescEntry}

%--------------------------------------------------------------------------------
\begin{iPageDescOpEntry}{Operation}
\begin{lstlisting}[style=iPageOpStyle]
RegR <- addOfs32( RegB, signExtend( ofs, 13 ));
RegR <- addOfs32( RegB, ( RegX ));
\end{lstlisting}
\end{iPageDescOpEntry}

%--------------------------------------------------------------------------------
\begin{iPageDescEntry}{Exceptions}
None.
\end{iPageDescEntry}

%--------------------------------------------------------------------------------
\begin{iPageDescEntry}{Notes}

The \texttt{LDO} instruction is often used to load immediate values. 

\begin{verbatim}
    LDI Rm, val -> LDO Rm, val( R0 )
\end{verbatim}

If value to load is larger than the 13-bit offset, a sequence of instruction is however necessary.

\end{iPageDescEntry}



